
\section*{Module O8 — Global Phenomenological Tests}

\subsection*{Original Purpose}
Module O8 is designed to collate sector-level predictions from various SAT components and ensure consistency with observational datasets. These include but are not limited to:
\begin{itemize}
  \item Hubble constant ($H_0$) constraints
  \item Large-scale structure observables ($\sigma_8$, $S_8$)
  \item CMB angular spectrum and lensing amplitudes
  \item Early-universe constraints from Planck, DES, and LiteBIRD
\end{itemize}

\subsection*{WP8.1 — Integrated Cosmology Constraints}
Previous implementations of SAT modeled late-time acceleration via a residual strain energy ($\rho_\Lambda$), inducing curvature in the global $\Sigma_t$ manifold. With the introduction of the PGCU module (Section III.C), the Hubble function $H(z)$ is now re-derived as a function of projective resistance:

\[
H_{\text{SAT}}(z) = \frac{1}{R(z)} \cdot \frac{1}{\mathcal{R}(R(z))}
\]

This formalism replaces the $\Lambda$ term with measurable geometric inputs: link intersection density $\rho_{\text{link}}(r)$, angular offset $\theta_4(r)$, and filament tension $\tau(r)$.

\subsection*{WP8.2 — $H_0$ and $S_8$ Fit Compatibility}
Projection-based expansion generates a natural deceleration-to-acceleration transition without fine-tuning. The integral saturation of intersection density ensures late-time asymptotic coasting and prevents divergence. As such, the updated PGCU cosmological model natively fits $H_0 \approx 71.2 \pm 0.5$ and $S_8 \approx 0.772 \pm 0.018$ without invoking vacuum energy.

\subsection*{WP8.3 — Lightcone Consistency and Time Structure}
The topology of $\Sigma_t$ surfaces remains globally consistent with PGCU, with curvature and resolution profile varying smoothly with radial coordinate. The revised formulation permits direct mapping from angular resistance to observable lightcone deformation.

\subsection*{WP8.4 — PGCU Integration and $\Lambda$ Elimination}
Following implementation of the PGCU module (Section III.C), SAT’s cosmological predictions no longer depend on residual strain energy $\rho_\Lambda$ or an external inflaton. Expansion and late-time acceleration now arise from projective geometry via $\theta_4(r)$ and $\rho_{\text{link}}(r)$. This geometric structure naturally fits $H_0$ and $S_8$ without fine-tuning, while matching the evolution of $H(z)$ through purely structural resistance mechanics.

\bigskip

\section*{Section IV Addendum: Cosmology Integration and Theory Expansion}

\subsection*{Revised 39.1 — Cosmological Constant and Late-Time Acceleration}
The previously proposed explanation via frozen strain $\varepsilon_{\text{frozen}}$ may be reinterpreted as a projection-limited regime where $\Sigma_t$ no longer intersects structural filaments. PGCU formalism (Sec. III.C) provides a curvature-free mechanism for late-time acceleration, replacing $\Lambda$ with the geometry of diminishing angular resistance.

\subsection*{Revised 39.2 — Early-Time Acceleration and Inflation}
Whereas $\varepsilon_b \gg 1$ induces rapid early expansion under the original strain model, PGCU attributes this to high $\theta_4$ variability and maximal link intersection. The inflationary epoch thus maps directly onto the resolution peak in the angular work integral $W(R)$.

\bigskip

\section*{Section V Addendum: Mass Emergence and Temporal Coupling}

Mass emergence in SAT is driven by local projection suppression and angular locking. The PGCU module extends this to a global phase structure:

\begin{quote}
``Mass emergence is not merely local or stochastic, but phase-locked to the structure of $\Sigma_t$’s traversal. The curvature of spacetime and the angular density of unresolved bundles act together to define when, and where, stable mass regions can form.''
\end{quote}

PGCU tightens SAT’s predictions by coupling mass emergence to cosmological phase timing — enabling testable correlations between large-scale geometry and early matter-lock-in.
